\section{Introduction}
\label{sec:intro}

Vision--language--action (VLA) policies are increasingly evaluated on \emph{spatial
generalization}: can a policy trained on objects at certain locations still act when the
object moves? The field's answer is a number---task success rate under object
displacement. That number is convenient, comparable, and, we argue, not diagnostic of
what it is taken to mean.

Consider a policy that, when the object is displaced, reaches toward the \emph{canonical}
location where it saw the object during training, rather than toward the object's new
position. If the task's grasp basin is wide---many pre-grasp poses lead to a successful
grasp---such an anchored reach can still succeed a nontrivial fraction of the time.
Success rate credits this as spatial generalization. A second policy that actually
\emph{grounds} to the displaced object and reaches for it produces the same success
number. Success rate cannot tell the two apart, yet they differ in exactly the capability
the metric claims to measure.

That VLA policies can anchor to training regularities is known, under names such as
causal confusion~\citep{dehaan2019causal} and ``memory traps''~\citep{afi2025memory}; that
success-only evaluation is imperfect has also been argued~\citep{liberox2026}. Our
contribution is neither. It is a \emph{readout} for a question these lines of work do not
answer: on the rollouts that \emph{succeed}, is the reach caused by the perceived object
position (grounding), or is it anchored to the canonical location and merely tolerated by
a wide basin?

We introduce an interventional, displacement-normalized, success-conditioned reach
readout. We displace the target object by $\dvec$ at episode start and measure where the
reach endpoint $\Ept$ lands, in the frame normalized by $\dvec$: $\uu = \langle \evec,
-\hat{\dvec}\rangle / \|\dvec\|$ runs from $0$ at the object to $1$ at the canonical
location, and $\rrho = \|\evec\|/\|\dvec\|$ is the normalized error, with $\evec = \Ept -
\Tpt$. We compute these at the pre-grasp, object-contact, and object-lift endpoints and,
crucially, split them by success and failure. The readout is cheap (it reuses the
rollout), model-agnostic (any policy that produces rollouts), and interpretable
($\uu \to 0$ means grounded, $\uu \to 1$ means anchored).

Using the readout on a competent SmolVLA~\citep{smolvla2025} across two LIBERO
suites~\citep{liu2023libero}, we surface a finding that success rate hides. The
\emph{same} counterfactual-augmentation fine-tuning improves displaced success in both
suites, but for opposite reasons (\Cref{fig:hero}). In wide-basin LIBERO-spatial,
successful reaches stay \emph{off the object} ($\usucc$ moves from $1.06$ to $0.51$):
success rises because the basin widens, not because the policy grounds. In tight-basin
LIBERO-object, successful reaches become \emph{object-near} ($\usucc$ moves from $0.67$ to
$0.21$; displaced success rises from $10\%$ to $27\%$): here the tight basin makes
grounding necessary for success, so the same intervention is forced to ground. Success
rate reports ``improvement'' in both and cannot distinguish the mechanisms.

\paragraph{Contributions.}
\begin{enumerate}[topsep=2pt,itemsep=1pt,leftmargin=1.4em]
\item A cheap, interventional, displacement-normalized, success-conditioned \emph{reach
readout} ($\uu$, $\rrho$; pre-grasp/contact/lift) that dissociates object-grounding from
grasp-basin tolerance, released as a drop-in evaluation tool (\Cref{sec:method}).
\item The \emph{basin-vs-grounding dissociation} on competent VLAs across two suites:
displaced successes can be anchored, which success rate cannot reveal (\Cref{sec:exp}).
\item The \emph{cross-regime finding}: the same counterfactual-fine-tuning success gain
reflects basin-widening in one suite and object-near reaching in another, invisible to
success rate (\Cref{sec:exp}).
\end{enumerate}
A boundary analysis (\Cref{sec:boundary}) shows the readout also \emph{rejects} tempting
mechanistic explanations that do not survive across suites, illustrating its
discriminating power. We scope our study to a single policy family and simulation;
\Cref{sec:conclusion} states these limits and the model-agnostic path forward.
